% =====================================================================
% MATERIAL SUPLEMENTAR — submetido como arquivo separado à RESS
% (Epidemiologia e Serviços de Saúde). NÃO é apêndice do corpo do
% artigo; não conta para o limite de 3.500 palavras / 5 figuras-tabelas.
% Acompanha: "Iniquidade regional no acesso à neuroimagem para a
% doença de Parkinson no Brasil, 2013–2025: estudo ecológico".
% Compilação: tectonic equipamentos-rm-parkinson-RESS-suplementar.tex
% =====================================================================
\documentclass[11pt,a4paper]{article}
\usepackage[T1]{fontenc}\usepackage[utf8]{inputenc}\usepackage[brazilian]{babel}
\usepackage{newtxtext}\usepackage{newtxmath}
\usepackage[a4paper,top=2cm,bottom=2cm,left=2cm,right=2cm]{geometry}
\usepackage{booktabs}\usepackage{array}\usepackage{caption}\usepackage{pdflscape}\usepackage{xcolor}
\usepackage[colorlinks=true,urlcolor=blue!55!black,linkcolor=black]{hyperref}
\usepackage{url}\usepackage{xurl}
\captionsetup{labelfont=bf,labelsep=endash,font=small,justification=raggedright,singlelinecheck=false}
\setlength{\parindent}{0pt}\setlength{\parskip}{6pt}
\begin{document}

{\large\bfseries Material suplementar}\\[2pt]
{\small Iniquidade regional no acesso à neuroimagem para a doença de
Parkinson no Brasil, 2013\textendash 2025: estudo ecológico de base
nacional.}

\vspace{6pt}
\textbf{Nota sobre proveniência dos dados.} Este material reúne o
painel completo por unidade federativa (UF) que, por limite de espaço,
não consta do corpo do artigo. Distinguem-se duas naturezas de dado:
(i) \emph{contagens brutas} da camada \emph{bronze} do
\textit{pipeline} \textemdash{} registros do CNES tal como ingeridos,
em modo \emph{append-only} (Tabela~S2); e (ii) \emph{grandezas
derivadas} (unidades-equipamento médias por ano e densidades per
capita) calculadas por agregação na camada \emph{gold}, deduplicadas
por estabelecimento (Tabela~S1). Densidades são, por definição,
agregados \textemdash{} não leituras diretas do dado bruto. Fontes:
CNES/DATASUS e estimativas populacionais do IBGE (2024). Código e dados
consolidados: \url{https://github.com/leonardochalhoub/mirante-dos-dados-br}.

\begin{landscape}
\begin{table}[htbp]\centering\footnotesize
\caption{Tabela S1 \textendash{} Parque instalado e densidade de
neuroimagem por unidade federativa, por modalidade (camada
\emph{gold}, valores agregados). Brasil, 2025. RM: ressonância
magnética; CT: tomografia computadorizada; PET/CT; câmara gama
(DAT-SPECT). Densidade por milhão de habitantes (Mhab); \%SUS:
participação do SUS nas unidades de RM.}
\setlength{\tabcolsep}{4pt}
\begin{tabular}{l r r r r r r r r r r}
\toprule
\textbf{UF} & \textbf{Pop.\ (mi)} & \textbf{RM} & \textbf{RM/Mhab} & \textbf{\%SUS} & \textbf{CT} & \textbf{CT/Mhab} & \textbf{PET/CT} & \textbf{PET/Mhab} & \textbf{Gama} & \textbf{Gama/Mhab} \\
\midrule
Acre (AC) & 0,88 & 5 & 5,7 & 60 & 22 & 24,9 & 0 & 0,0 & 31 & 35,1 \\
Alagoas (AL) & 3,22 & 38 & 11,8 & 61 & 74 & 23,0 & 3 & 0,9 & 465 & 144,4 \\
Amazonas (AM) & 4,32 & 37 & 8,6 & 76 & 79 & 18,3 & 1 & 0,2 & 299 & 69,2 \\
Amapá (AP) & 0,81 & 14 & 17,4 & 71 & 29 & 36,0 & 0 & 0,0 & 10 & 12,4 \\
Bahia (BA) & 14,87 & 217 & 14,6 & 57 & 390 & 26,2 & 7 & 0,5 & 1.426 & 95,9 \\
Ceará (CE) & 9,27 & 92 & 9,9 & 37 & 269 & 29,0 & 2 & 0,2 & 1.019 & 109,9 \\
Distrito Federal (DF) & 3,00 & 138 & 46,0 & 30 & 256 & 85,4 & 10 & 3,3 & 102 & 34,0 \\
Espírito Santo (ES) & 4,13 & 76 & 18,4 & 45 & 140 & 33,9 & 1 & 0,2 & 212 & 51,4 \\
Goiás (GO) & 7,42 & 144 & 19,4 & 63 & 356 & 48,0 & 2 & 0,3 & 259 & 34,9 \\
Maranhão (MA) & 7,02 & 85 & 12,1 & 51 & 230 & 32,8 & 4 & 0,6 & 803 & 114,4 \\
Minas Gerais (MG) & 21,39 & 375 & 17,5 & 50 & 886 & 41,4 & 24 & 1,1 & 1.560 & 72,9 \\
Mato Grosso do Sul (MS) & 2,92 & 57 & 19,5 & 37 & 142 & 48,6 & 2 & 0,7 & 329 & 112,5 \\
Mato Grosso (MT) & 3,89 & 94 & 24,1 & 63 & 202 & 51,9 & 2 & 0,5 & 138 & 35,4 \\
Pará (PA) & 8,71 & 108 & 12,4 & 53 & 252 & 28,9 & 3 & 0,3 & 405 & 46,5 \\
Paraíba (PB) & 4,16 & 58 & 13,9 & 55 & 150 & 36,0 & 4 & 1,0 & 1.198 & 287,7 \\
Pernambuco (PE) & 9,56 & 131 & 13,7 & 55 & 212 & 22,2 & 9 & 0,9 & 579 & 60,6 \\
Piauí (PI) & 3,38 & 36 & 10,6 & 67 & 123 & 36,3 & 1 & 0,3 & 571 & 168,7 \\
Paraná (PR) & 11,89 & 236 & 19,8 & 50 & 464 & 39,0 & 10 & 0,8 & 736 & 61,9 \\
Rio de Janeiro (RJ) & 17,22 & 466 & 27,1 & 31 & 894 & 51,9 & 23 & 1,3 & 1.422 & 82,6 \\
Rio Grande do Norte (RN) & 3,46 & 36 & 10,4 & 83 & 68 & 19,7 & 2 & 0,6 & 233 & 67,4 \\
Rondônia (RO) & 1,75 & 53 & 30,3 & 66 & 116 & 66,2 & 0 & 0,0 & 117 & 66,8 \\
Roraima (RR) & 0,74 & 6 & 8,1 & 100 & 15 & 20,3 & 0 & 0,0 & 35 & 47,4 \\
Rio Grande do Sul (RS) & 11,23 & 264 & 23,5 & 52 & 507 & 45,1 & 7 & 0,6 & 323 & 28,8 \\
Santa Catarina (SC) & 8,19 & 171 & 20,9 & 69 & 365 & 44,6 & 5 & 0,6 & 867 & 105,9 \\
Sergipe (SE) & 2,30 & 27 & 11,7 & 41 & 44 & 19,1 & 7 & 3,0 & 82 & 35,7 \\
São Paulo (SP) & 46,08 & 896 & 19,4 & 39 & 1.657 & 36,0 & 37 & 0,8 & 2.809 & 61,0 \\
Tocantins (TO) & 1,59 & 38 & 23,9 & 71 & 59 & 37,2 & 0 & 0,0 & 60 & 37,8 \\
\midrule
\textbf{Brasil} & \textbf{213,4} & \textbf{3.898} & \textbf{18,3} & \textbf{48} & \textbf{8.001} & \textbf{37,5} & \textbf{166} & \textbf{0,8} & \textbf{16.090} & \textbf{75,4} \\
\bottomrule
\end{tabular}
\end{table}
\end{landscape}

\begin{table}[htbp]\centering\small
\caption{Tabela S2 \textendash{} Volume bruto da camada \emph{bronze}
por unidade federativa e modalidade, 2013\textendash 2025 (registros
CNES \emph{append-only} e estabelecimentos distintos com RM). Reflete o
dado-fonte antes de deduplicação/agregação; não corresponde ao número
de aparelhos.}
\setlength{\tabcolsep}{6pt}
\begin{tabular}{l r r r r r}
\toprule
\textbf{UF} & \textbf{RM (reg.)} & \textbf{RM (estab.)} & \textbf{CT (reg.)} & \textbf{PET/CT (reg.)} & \textbf{Gama (reg.)} \\
\midrule
Acre (AC) & 946 & 5 & 2.876 & 0 & 1.694 \\
Alagoas (AL) & 5.310 & 32 & 13.562 & 214 & 14.820 \\
Amazonas (AM) & 5.216 & 36 & 11.946 & 130 & 3.494 \\
Amapá (AP) & 1.482 & 14 & 3.852 & 0 & 734 \\
Bahia (BA) & 36.964 & 212 & 67.208 & 506 & 37.604 \\
Ceará (CE) & 15.936 & 96 & 46.572 & 82 & 19.348 \\
Distrito Federal (DF) & 20.998 & 152 & 35.938 & 1.246 & 6.832 \\
Espírito Santo (ES) & 15.942 & 75 & 26.952 & 92 & 8.514 \\
Goiás (GO) & 22.910 & 141 & 61.548 & 184 & 11.622 \\
Maranhão (MA) & 10.246 & 88 & 35.356 & 722 & 21.472 \\
Minas Gerais (MG) & 69.650 & 380 & 155.182 & 1.890 & 31.442 \\
Mato Grosso do Sul (MS) & 9.866 & 61 & 21.908 & 320 & 4.834 \\
Mato Grosso (MT) & 14.718 & 109 & 34.092 & 348 & 6.194 \\
Pará (PA) & 17.778 & 115 & 41.200 & 640 & 12.512 \\
Paraíba (PB) & 9.860 & 55 & 26.460 & 540 & 37.452 \\
Pernambuco (PE) & 20.834 & 116 & 35.348 & 1.096 & 21.234 \\
Piauí (PI) & 5.818 & 35 & 20.014 & 240 & 14.478 \\
Paraná (PR) & 42.096 & 225 & 89.260 & 1.242 & 14.376 \\
Rio de Janeiro (RJ) & 84.778 & 550 & 158.070 & 2.630 & 31.018 \\
Rio Grande do Norte (RN) & 6.362 & 39 & 13.744 & 480 & 6.070 \\
Rondônia (RO) & 7.492 & 54 & 16.542 & 8 & 3.454 \\
Roraima (RR) & 1.048 & 6 & 2.370 & 0 & 1.418 \\
Rio Grande do Sul (RS) & 50.652 & 248 & 97.314 & 1.390 & 15.810 \\
Santa Catarina (SC) & 29.760 & 165 & 59.916 & 858 & 21.448 \\
Sergipe (SE) & 3.382 & 22 & 8.192 & 562 & 3.698 \\
São Paulo (SP) & 156.712 & 834 & 313.508 & 5.528 & 64.064 \\
Tocantins (TO) & 5.794 & 37 & 9.588 & 0 & 2.984 \\
\midrule
\textbf{Brasil} & \textbf{672.550} & \textbf{3.902} & \textbf{1.408.518} & \textbf{20.948} & \textbf{418.620} \\
\bottomrule
\end{tabular}
\par\vspace{3pt}{\footnotesize\raggedright ``reg.'' = registros brutos
bronze (append-only, 2013\textendash 2025); ``estab.'' = CNES distintos
com ao menos um registro de RM no período. O total nacional bruto das
quatro modalidades é 2.520.636 registros em 17.846 estabelecimentos
distintos.\par}
\end{table}

\end{document}
